\pdfoutput=1
\documentclass[a4paper,12pt]{article}

\usepackage{fix-cm}
\usepackage{amsfonts}
\usepackage{mathrsfs}
\usepackage{amsmath}
\usepackage{amssymb}
\usepackage{framed}

\usepackage[medium]{titlesec}
\usepackage{bm}
\usepackage{cite}

\usepackage[normalem]{ulem}
\usepackage{extarrows}
\usepackage{slashed}
\usepackage{isodateo}
\usepackage{graphicx}
\usepackage[dvipsnames]{xcolor}
\usepackage[bookmarksnumbered=true,bookmarksopen=true]{hyperref}
 \hypersetup{colorlinks,%
             linkcolor=NavyBlue, %
             citecolor=PineGreen, %
             urlcolor=PineGreen}
\usepackage[hmargin=.7in,vmargin=1.1in]{geometry}
\usepackage{indentfirst}
\usepackage{booktabs}

\usepackage{bbm}

\linespread{1.1}

\newcommand{\FR}[2]{\displaystyle\frac{\,{#1}\,}{#2}}
\newcommand{\fr}[2]{\mbox{$\frac{\,{#1}\,}{#2}$}}
\newcommand{\n}{\nonumber}
\renewcommand{\rm}{\mathrm}
\renewcommand{\thefootnote}{\arabic{footnote}}

\graphicspath{{fig/}}


\newcommand{\red}{\color{red}}
\newcommand{\blue}{\color{blue}}
\newcommand{\gray}{\color{gray}}

\def\bge{\begin{equation}}
\def\ede{\end{equation}}
\def\bga{\begin{aligned}}
\def\eda{\end{aligned}}
\def\ie{\begin{equation}\begin{aligned}}
\def\fe{\end{aligned}\end{equation}}
\def\bgb{\begin{bmatrix}}
\def\edb{\end{bmatrix}}
\def\bgp{\begin{pmatrix}}
\def\edp{\end{pmatrix}}
\def\bgm{\begin{matrix}}
\def\edm{\end{matrix}}
\def\bgs{\begin{subequations}}
\def\eds{\end{subequations}}
\newcommand{\order}[1]{\mathcal{O}({#1})}
\def\di{{\mathrm{d}}}
\def\D{{\mathrm{D}}}
\def\Di{{\mathcal{D}}}

\def\T{\mathcal{T}}
\def\mb{\mathbf}
\def\ms{\mathscr}
\def\mf{\mathfrak}
\def\mc{\mathcal}
\def\cl{\mathrm{cl}}
\def\pd{\partial}
\def\ld{{\mathscr{L}}}
\def\hd{{\mathscr{H}}}
\def\z{{\bar{z}}}
\def\la{\langle}\def\ra{\rangle}
\def\sla{\slashed}
\def\const{\mathrm{const.}}
\setlength\unitlength{1mm}
\def\tr{\mathrm{\,tr\,}}
\def\Tr{\mathrm{\,Tr\,}}
\def\Det{\mathrm{\,Det\,}}


\def\To{\Rightarrow}
\def\ii{\mathrm{i}}

\def\al{\alpha}
\def\be{\beta}
\def\ga{\gamma}
\def\de{\delta}
\def\ep{\epsilon}
\def\ka{\kappa}
\def\lam{\lambda}
\def\rh{\rho}
\def\si{\sigma}
\def\ze{\zeta}

\def\Wh{\mathrm{W}}


\def\aa{\mathsf{a}}
\def\bb{\mathsf{b}}
\def\cc{\mathsf{c}}
\def\dd{\mathsf{d}}

\def\C{\mathbb{C}}
\def\R{\mathbb{R}}

\newcommand{\mN}{\mathcal{N}}
\newcommand{\mG}{\mathcal{G}}
\newcommand{\mK}{\mathcal{K}}
\def\to{\rightarrow}

\def\Mp{M_{\text{Pl}}}


\def\Re{\mathrm{Re}\,}
\def\Im{\mathrm{Im}\,}

\def\ad{\mathrm{\,ad\,}}

\def\2F1{{}_2\mathrm{F}_1}
\def\3F2{{}_3\mathrm{F}_2}

\def\nut{\wt{\nu}}
%\newcommand{\nut}{\widetilde{\nu}}


\def\tf{{\tau_0}}

\usepackage{mdframed} 
                   

\newmdenv[skipabove=0pt,%
          skipbelow=5pt,%
          leftmargin=0pt,%
          rightmargin=0pt,%
          innertopmargin=-5pt,%
          innerbottommargin=7pt,%
          innerleftmargin=2pt,%
          innerrightmargin=2pt,%
          splittopskip=0pt,%
          splitbottomskip=0pt,%
          linewidth=0pt,%
          nobreak=true]%
          {keyeqn2}                   
                   
                   
\newmdenv[backgroundcolor=gray!15,%
          skipabove=0pt,%
          skipbelow=5pt,%
          leftmargin=0pt,%
          rightmargin=0pt,%
          innertopmargin=-5pt,%
          innerbottommargin=7pt,%
          innerleftmargin=2pt,%
          innerrightmargin=2pt,%
          splittopskip=0pt,%
          splitbottomskip=0pt,%
          linewidth=0pt,%
          nobreak=true]%
          {keyeqn}
          
\usepackage{titlesec}          
\titleformat{\section}
{\normalfont\fontsize{15}{20}\bfseries}{\thesection}{1em}{}

\newcommand{\ob}[1]{\mkern 2mu \overline{\mkern -2mu #1 \mkern -2mu}\mkern 2mu}
\newcommand{\wt}[1]{\mkern 2mu \widetilde{\mkern -2mu #1 \mkern -2mu}\mkern 2mu}
\newcommand{\wh}[1]{\mkern 2mu \widehat{\mkern-2mu#1\mkern-2mu}\mkern 2mu}

\newcommand{\fnemail}[1]{\footnote{Email: \href{mailto:#1}{\nolinkurl{#1}}}}


\newcommand{\zq}[1]{{\color[rgb]{.1,.6,.3}{[{ZQ:} #1]}}}
\newcommand{\cjq}[1]{{\color[rgb]{.1,.6,.3}{[{CJQ:} #1]}}}
\newcommand{\tyx}[1]{{\color[rgb]{.1,.6,.3}{[{YT:} #1]}}}
\begin{document}

\title{\Large\textbf{IBP for V++ Massive Bubble\\[2mm]}}

\author{Zhehan Qin$^{\,a,b\,}$\fnemail{qzh21@mails.tsinghua.edu.cn}\\[5mm]
	$^a\,$\normalsize{\emph{Department of Physics, Tsinghua University, Beijing 100084, China} }\\
	$^b\,$\normalsize{\emph{Department of Applied Mathematics and Theoretical Physics, University of Cambridge,}}\\\normalsize{\emph{Wilberforce Road, Cambridge, CB3 0WA, UK}}\\
}

\date{}
\maketitle

\vspace{20mm}

\begin{abstract}
In this note we construct IBP relations for the one-loop two-point (bubble) diagram in de Sitter space with both vertices being plus (V++) and internal propagators massive. We follow the structure of the reference note for V+- massive bubble, adding only the modifications required by the $\theta$-function time-ordering. The momentum IBP relations remain unchanged, while time IBP relations acquire tadpole terms from derivatives of $\theta$-functions. We provide the complete IBP system ready for Mathematica implementation.
	\vspace{10mm}
\end{abstract}

\newpage
\tableofcontents


\section{Background, goals, and structure}
The goal is to construct practical reduction relations for the massive V++ de Sitter bubble, including the tadpole families induced by time ordering. Differential operators in external kinematics map integrals within the same families; after reduction to a master basis, these operators induce a first-order linear system for the master-integral vector.

We organize the note as follows. Section~2 fixes conventions and defines the $B^{(12)}$, $B^{(21)}$ and tadpole $R$ families. Section~3 lists the IBP system and the kinematic differential operators. Section~4 records Baikov and dimension-shift relations for the momentum integral. Section~5 summarizes a $d\log$/intersection-theory viewpoint tailored to the present special-function setting.

\section{Basic definitions, notation, and integral families}
\label{sec:definitions}
\paragraph{Notations} The metric reads: $\di s^2=a^2(-\di\tau^2+\di \mb x^2)$ with scale factor $a(\tau)=-1/(H\tau)$ and $\tau\in(-\infty,0)$. We set $H\equiv 1$ throughout. We use $\theta_{ij}\equiv\theta(\tau_i-\tau_j)$ and $\theta(0)=1/2$. We work in $D$ spatial dimensions (so the loop momentum integration is $\int\di^D\mb q$); in dimensional regularization one may take $D=3-2\ep$. For simplicity we take all mass parameters equal, $\nu_1=\nu_2=\nu_3=\nu_4\equiv\nu$, and we also specialize to equal external energies $P_1=P_2$.

\paragraph{Hankel building blocks} Following the notation of arXiv:2401.00129v5 (Eqs.~3.11--3.12), we define the rescaled Hankel functions
\bge
h_\nu^{(1)}(x) \equiv x^{-\nu}\mathrm H_{
\nu}^{(1)}(x),\qquad h_\nu^{(2)}(x) \equiv x^{-\nu}\mathrm H_{
\nu^*}^{(2)}(x),\qquad 
\nu\in\mathbb C,\quad x>0.
\ede
They satisfy (Eq.~(3.13) in arXiv:2401.00129v5)
\begin{equation}
	\label{eq:hdef}
	h''(x) + \FR{2\nu+1}{x}h'(x) + h(x) = 0.
\end{equation}
The mass parameter is $m^2=9/4-\nu^2$. We also define derivatives
\bge
h_{\nu,0}(x)\equiv h_\nu(x),\qquad h_{\nu,n} (x)  \equiv  \FR{\di^n}{\di x^n}h(x),
\ede
For the two branches we also write
\bge
h^{(1)}_{\nu,n}(x)\equiv \FR{\di^n}{\di x^n}h^{(1)}_\nu(x),\qquad
h^{(2)}_{\nu,n}(x)\equiv \FR{\di^n}{\di x^n}h^{(2)}_\nu(x).
\ede
When the mass label is clear from context, we write $h_n\equiv h_{\nu,n}$ and similarly $h^{(1,2)}_n\equiv h^{(1,2)}_{\nu,n}$.
In this note we consider a single mass scale, so we set all $\nu_i=\nu$ everywhere.

\paragraph{V++ bulk-to-bulk propagator and time ordering} For a massive scalar on the $+$ branch, the Schwinger--Keldysh propagator contains two time orderings,
\bge
B^{++}_\nu(k;\tau_1,\tau_2)\;\propto\;
h_0^{(1)}(-k\tau_1)\,h_0^{(2)}(-k\tau_2)\,\theta_{12}
\;+\;
h_0^{(2)}(-k\tau_1)\,h_0^{(1)}(-k\tau_2)\,\theta_{21},
\ede
with $\tau_1,\tau_2\in(-\infty,0)$. For the bubble, the product of two $B^{++}$ propagators naturally splits into two contributions proportional to $\theta_{12}$ and $\theta_{21}$ respectively. This motivates treating the two orderings as two integral families related by symmetry.

\paragraph{Two integral families} To clearly distinguish the two time-orderings, we define two integral families with explicit $\theta$-function labels:

\paragraph{Family with $\theta_{12}$}
\begin{align}
	B^{(12)}[\{n_1,n_2,n_3,n_4\},\{a_1,a_2\},\{b_1,b_2\}] \equiv & \int_{-\infty}^0 \di\tau_1\di\tau_2 \int \FR{\di^D \mb q}{(2\pi)^D}\, e^{-\ii P_1\tau_1 - \ii P_2 \tau_2}\,
	\frac{(-\tau_1)^{a_1}(-\tau_2)^{a_2}}{|\mb q|^{b_1}|\mb k_s+\mb q|^{b_2}} \nonumber \\
	& \times h_{n_1}^{(1)}(-q\tau_1)h_{n_2}^{(2)}(-q\tau_2)\nonumber\\
	& \times h_{n_3}^{(1)}(-|\mb q+\mb k_s|\tau_1)h_{n_4}^{(2)}(-|\mb q+\mb k_s|\tau_2)\,\theta_{12}.
\end{align}

\paragraph{Family with $\theta_{21}$}
\begin{align}
	B^{(21)}[\{n_1,n_2,n_3,n_4\},\{a_1,a_2\},\{b_1,b_2\}] \equiv & \int_{-\infty}^0 \di\tau_1\di\tau_2 \int \FR{\di^D \mb q}{(2\pi)^D}\, e^{-\ii P_1\tau_1 - \ii P_2 \tau_2}\,
	\frac{(-\tau_1)^{a_1}(-\tau_2)^{a_2}}{|\mb q|^{b_1}|\mb k_s+\mb q|^{b_2}} \nonumber \\
	& \times h_{n_1}^{(2)}(-q\tau_1)h_{n_2}^{(1)}(-q\tau_2)\nonumber\\
	& \times h_{n_3}^{(2)}(-|\mb q+\mb k_s|\tau_1)h_{n_4}^{(1)}(-|\mb q+\mb k_s|\tau_2)\,\theta_{21}.
\end{align}

\paragraph{Relation between the two families}
The two families are related by exchanging time variables and external energies:
\begin{align}
	B^{(21)}[\{n_1,n_2,n_3,n_4\},\{a_1,a_2\},\{b_1,b_2\}; P_1,P_2] = B^{(12)}[\{n_2,n_1,n_4,n_3\},\{a_2,a_1\},\{b_1,b_2\}; P_2,P_1].
\end{align}

\paragraph{Complete V++ correlator}
The physical V++ bubble correlator is the sum:
\begin{align}
	\mathcal{V}_{++} = B^{(12)} + B^{(21)}.
\end{align}

\paragraph{Tadpole families} Time IBP relations involve total derivatives $\int \pd_{\tau_i}[\cdots]=0$. When the integrand contains $\theta_{12}$ or $\theta_{21}$, one generates coincidence terms proportional to $\delta(\tau_1-\tau_2)$, which reduce the two-time integral to a one-time integral. We treat these coincidence contributions as tadpole families $R_1,R_2$, which depend only on $P_0\equiv P_1+P_2$ (and we keep $P_0$ implicit in the notation).
\begin{align}
	R_1[\{n_3,n_4\},\{a\},\{b_1,b_2\}] & \equiv e^{\pi \Im [\nu]}\frac{4i}{\pi}\int_{-\infty}^0 \di\tau \int \FR{\di^D \mb q}{(2\pi)^D}\, e^{-\ii P_0\tau} \frac{(-\tau)^{a}}{|\mb q|^{b_1}|\mb k_s+\mb q|^{b_2}} \nonumber \\
	                                        & \quad\times h_{n_3}^{(1)}(-|\mb q+\mb k_s|\tau)h_{n_4}^{(2)}(-|\mb q+\mb k_s|\tau),                                                                                                                 \\
	R_2[\{n_1,n_2\},\{a\},\{b_1,b_2\}] & \equiv e^{\pi \Im [\nu]}\frac{4i}{\pi}\int_{-\infty}^0 \di\tau \int \FR{\di^D \mb q}{(2\pi)^D}\, e^{-\ii P_0\tau} \frac{(-\tau)^{a}}{|\mb q|^{b_1}|\mb k_s+\mb q|^{b_2}} \nonumber \\
	                                        & \quad\times h_{n_1}^{(1)}(-q\tau)h_{n_2}^{(2)}(-q\tau).
\end{align}

\paragraph{Direct Tadpole Family}
According to the Feynman rules, the tadpole diagram is given by the one-time integral of the equal-time propagator. Adopting the regularization $\theta(0)=1/2$, the equal-time V++ propagator is
\begin{equation}
B^{++}_
u(q;\tau,\tau) = \frac{1}{2}\left[h^{(1)}_
u(-q\tau)h^{(2)}_
u(-q\tau) + h^{(2)}_
u(-q\tau)h^{(1)}_
u(-q\tau)\right] = h^{(1)}_
u(-q\tau)h^{(2)}_
u(-q\tau).
\end{equation}
We define the direct tadpole family $T$ corresponding to this diagrammatic contribution:
\begin{align}
	T[\{n_1,n_2\},\{a\},\{b\}] & \equiv \int_{-\infty}^0 \di\tau \int \FR{\di^D \mb q}{(2\pi)^D}\, e^{-\ii P_0\tau} \frac{(-\tau)^{a}}{|\mb q|^{b}} \nonumber \\
	                                        & \quad\times h_{n_1}^{(1)}(-q\tau)h_{n_2}^{(2)}(-q\tau).
\end{align}
The families $R_1,R_2$ instead arise as remaining terms in time IBP for the bubble: one propagator is contracted using the Wronskian, producing the universal prefactor $e^{\pi \Im [\nu]}\frac{4i}{\pi}$ and an extra $(-k\tau)^{-2
u-1}$ factor for the contracted line. In the integral-family language this is implemented as shifts of the $\tau$-power and the corresponding $|k|$-power indices (e.g. $a\to a-2
u-1$ and $b\to b+2
u+1$), together with the constraint $\delta_{n_i+n_j,1}$ from the contraction formula.

\section{Reduction relations for the one-loop two-point bubble}
\label{sec:reductions}
\subsection{Top-sector families \texorpdfstring{$B^{(12)},B^{(21)}$}{B(12), B(21)}}
\paragraph{Origin of the tadpole (remaining) terms}
Time IBP relations involve total derivatives $\int \pd_{\tau_i}[\cdots]=0$. With $\theta_{12}$ or $\theta_{21}$ present, derivatives generate coincidence terms proportional to $\delta(\tau_1-\tau_2)$. For Hankel-type integrands, the coincidence terms collapse one time integration using the Wronskian of the $h^{(1,2)}_\nu$ pair, producing one-time tadpole families. In the notation of arXiv:2401.00129v5 (Section~3.7), in the time derivation of   $B_{++}$,  the $\theta$ functions yield a remaining term 
\bge
h_{n_1}^{(1)}(-k\tau_i)\,h_{n_2}^{(2)}(-k\tau_i) - h_{n_1}^{(2)}(-k\tau_i)\,h_{n_2}^{(1)}(-k\tau_i)
= \delta_{n_1+ n_2,1}e^{ \pi \text{Im}[\nu]} \frac{  4 i }{\pi}  (-k\tau_i)^{-2\nu-1} \,  .  \label{eq_propcontract}
\ede
which, in our integral-family language, is implemented as a shift of the corresponding $|k|$-power index by $2\nu+1$ and the $\tau$-power index by $-2\nu-1$,  together with the contraction of the two-time integrations by one integral in the remaining family.

In particular, since the $R$ families originate from contracting one propagator in the top sector, the contracted line carries an extra $(-k)^{-2\nu-1}$ factor from the Wronskian. In the Mathematica/Kira implementation we keep the momentum baselines trivial,
\[
|q|^{-b_1}\,|\bm k_s+\bm q|^{-b_2},
\qquad b_{10}=0, \quad b_{20}=-2\nu,
\]
and the remaining-term factor is tracked explicitly by the $b$ shifts (e.g. $b_1\to b_1+2\nu+1$) in Eqs.~(\ref{eq:timeIBP_tau1})--(\ref{eq:timeIBP_tau2}).


\paragraph{Time IBP with $\theta$-function tadpole terms}
In the following, $B$ denotes $B^{(12)}+B^{(21)}$. The combination of the two parts can employ the propagator contraction formula, thereby simplifying the calculation. Momentum IBP and equation-of-motion relations apply to the families first. The only new ingredient in V++ (compare to V+-) is the $\theta$ function contributions in the time IBP.

For $\pd_{\tau_1}$:
\begin{align}
	0 = & - \ii P_1 B[\{n_1,n_2,n_3,n_4\},\{a_1,a_2\},\{b_1,b_2\}] - a_1B[\{n_1,n_2,n_3,n_4\},\{a_1-1,a_2\},\{b_1,b_2\}]\nonumber \\
	    & - B[\{n_1+1,n_2,n_3,n_4\},\{a_1,a_2\},\{b_1-1,b_2\}] - B[\{n_1,n_2,n_3+1,n_4\},\{a_1,a_2\},\{b_1,b_2-1\}]\nonumber      \\
	    & +\delta_{n_1+n_2,1}(-1)^{n_1+1}R_1[\{n_3,n_4\},\{a_1+a_2-2\nu-1\},\{b_1+2\nu+1,b_2\}]\nonumber                \\
	    & +\delta_{n_3+n_4,1}(-1)^{n_3+1}R_2[\{n_1,n_2\},\{a_1+a_2-2\nu-1\},\{b_1,b_2+2\nu+1\}], \label{eq:timeIBP_tau1}
\end{align}
and for $\partial_{\tau_2}$:
\begin{align}
	0 = & - \ii P_2 B[\{n_1,n_2,n_3,n_4\},\{a_1,a_2\},\{b_1,b_2\}] - a_2B[\{n_1,n_2,n_3,n_4\},\{a_1,a_2-1\},\{b_1,b_2\}]\nonumber \\
	    & - B[\{n_1,n_2+1,n_3,n_4\},\{a_1,a_2\},\{b_1-1,b_2\}] - B[\{n_1,n_2,n_3,n_4+1\},\{a_1,a_2\},\{b_1,b_2-1\}]\nonumber      \\
	    & +\delta_{n_1+n_2,1}(-1)^{n_2+1}R_1[\{n_3,n_4\},\{a_1+a_2-2\nu-1\},\{b_1+2\nu+1,b_2\}]\nonumber                \\
	    & +\delta_{n_3+n_4,1}(-1)^{n_4+1}R_2[\{n_1,n_2\},\{a_1+a_2-2\nu-1\},\{b_1,b_2+2\nu+1\}], \label{eq:timeIBP_tau2}
\end{align}
where $\nu$ is the common mass parameter.

The tadpole families $R_1,R_2$ are defined in Section~\ref{sec:definitions}.

\paragraph{Momentum IBP}
Momentum IBP relations are unchanged compared to the V$+-$ bubble, because $\pd_{\bm q^i}\theta_{ij}=0$. The derivation can be organized in the same chain-rule form as in Section~1.6 of the reference note by introducing
\begin{align}
	x_1 \equiv |\bm q|,\qquad x_2 \equiv |\bm q+\bm k_s|,\qquad
	\mathcal{O}_1 \equiv \bm q^i \pd_{\bm q^i},\qquad
	\mathcal{O}_2 \equiv (\bm q^i+\bm k_s^i)\pd_{\bm q^i},
\end{align}
whose actions on $(x_1,x_2)$ are
\begin{align}
	\mathcal{O}_1 x_1 & = x_1, \nonumber\\
	\mathcal{O}_1 x_2 & = \frac{1}{2}x_2^{-1}(x_1^2 + x_2^2 - k_s^2), \nonumber\\
	\mathcal{O}_2 x_1 & = \frac{1}{2}x_1^{-1}(x_1^2 + x_2^2 - k_s^2), \nonumber\\
	\mathcal{O}_2 x_2 & = x_2.
\end{align}
Equivalently, for any function $F(x_1,x_2)$ one has the chain rule
\begin{align}
	\mathcal{O}_\alpha F = (\mathcal{O}_\alpha x_1)\,\partial_{x_1}F + (\mathcal{O}_\alpha x_2)\,\partial_{x_2}F,\qquad \alpha=1,2.
\end{align}

From $\int \pd_{\bm q^i}[\bm q^i \cdots ]=0$ one may write an intermediate step in terms of the composite operator acting on the full integrand,
\begin{align}
	0 &= \int_{-\infty}^0 \di\tau_1\di\tau_2\int \FR{\di^D \bm q}{(2\pi)^D}\,
	\Big[D+\mathcal{O}_1\Big] \nonumber\\ 
	& \Bigg[
	e^{-\ii P_1\tau_1-\ii P_2\tau_2}
	\frac{(-\tau_1)^{a_1}(-\tau_2)^{a_2}}{x_1^{b_1}x_2^{b_2}}\,
	h_{n_1}^{(1)}(-x_1\tau_1)h_{n_2}^{(2)}(-x_1\tau_2) \times h_{n_3}^{(1)}(-x_2\tau_1)h_{n_4}^{(2)}(-x_2\tau_2)\Bigg] \nonumber\\
	& =  \int_{-\infty}^0 \di\tau_1\di\tau_2\int \FR{\di^D \bm q}{(2\pi)^D}\,
	\Bigg[D+(\mathcal{O}_1 x_1)\partial_{x_1} +(\mathcal{O}_1 x_2)\partial_{x_2}\Bigg] \nonumber\\
	&\Bigg[
	e^{-\ii P_1\tau_1-\ii P_2\tau_2}
	\frac{(-\tau_1)^{a_1}(-\tau_2)^{a_2}}{x_1^{b_1}x_2^{b_2}}\,
	h_{n_1}^{(1)}(-x_1\tau_1)h_{n_2}^{(2)}(-x_1\tau_2)
	  \times h_{n_3}^{(1)}(-x_2\tau_1)h_{n_4}^{(2)}(-x_2\tau_2)\Bigg]. \label{eq:momIBP_qdot_chainrule}
\end{align}
Expanding the derivatives using $\partial_x x^{-b}=-(b/x)x^{-b}$ and $\partial_x h_n(-x\tau)=(-\tau)\,h_{n+1}(-x\tau)$ yields the closed relation in terms of shifted $B$ integrals:
\begin{align}
	0 ={}& D\,B[\{n_1,n_2,n_3,n_4\},\{a_1,a_2\},\{b_1,b_2\}] \nonumber\\
	&- b_1\,B[\{n_1,n_2,n_3,n_4\},\{a_1,a_2\},\{b_1,b_2\}] \nonumber \\
	&- \FR12 b_2\Big(
	B[\{n_1,n_2,n_3,n_4\},\{a_1,a_2\},\{b_1,b_2\}] \nonumber\\
	&\qquad + B[\{n_1,n_2,n_3,n_4\},\{a_1,a_2\},\{b_1-2,b_2+2\}] \nonumber\\
	&\qquad -k_s^2\,B[\{n_1,n_2,n_3,n_4\},\{a_1,a_2\},\{b_1,b_2+2\}]
	\Big) \nonumber \\
	&+B[\{n_1+1,n_2,n_3,n_4\},\{a_1+1,a_2\},\{b_1-1,b_2\}] \nonumber\\
	&+ B[\{n_1,n_2+1,n_3,n_4\},\{a_1,a_2+1\},\{b_1-1,b_2\}] \nonumber \\
	&+\FR12 \Big(
	B[\{n_1,n_2,n_3+1,n_4\},\{a_1+1,a_2\},\{b_1,b_2-1\}] \nonumber\\
	&\qquad + B[\{n_1,n_2,n_3+1,n_4\},\{a_1+1,a_2\},\{b_1-2,b_2+1\}] \nonumber\\
	&\qquad -k_s^2\,B[\{n_1,n_2,n_3+1,n_4\},\{a_1+1,a_2\},\{b_1,b_2+1\}]
	\Big) \nonumber \\
	&+\FR12\Big(
	B[\{n_1,n_2,n_3,n_4+1\},\{a_1,a_2+1\},\{b_1,b_2-1\}] \nonumber\\
	&\qquad + B[\{n_1,n_2,n_3,n_4+1\},\{a_1,a_2+1\},\{b_1-2,b_2+1\}] \nonumber\\
	&\qquad -k_s^2\,B[\{n_1,n_2,n_3,n_4+1\},\{a_1,a_2+1\},\{b_1,b_2+1\}]
	\Big),
\label{eq:momIBP_qdot}
\end{align}
where $D$ is the spatial dimension (in dimensional regularization $D=3-2\ep$). The second momentum IBP can be written in the same chain-rule form as
\begin{align}
	0 ={}& \int_{-\infty}^0 \di\tau_1\di\tau_2\int \FR{\di^D \bm q}{(2\pi)^D}\,
	\Big[D+\mathcal{O}_2\Big]\nonumber\\
	&\Bigg[
	e^{-\ii P_1\tau_1-\ii P_2\tau_2}
	\frac{(-\tau_1)^{a_1}(-\tau_2)^{a_2}}{x_1^{b_1}x_2^{b_2}}\,
	h_{n_1}^{(1)}(-x_1\tau_1)h_{n_2}^{(2)}(-x_1\tau_2) \times h_{n_3}^{(1)}(-x_2\tau_1)h_{n_4}^{(2)}(-x_2\tau_2)\Bigg], \label{eq:momIBP_qk_chainrule}
\end{align}
and, after expansion, it gives the symmetric relation with the exchanges
\begin{align}
	\{n_1,n_2,n_3,n_4\} &\leftrightarrow \{n_3,n_4,n_1,n_2\}, \nonumber\\
	\{b_1,b_2\} &\leftrightarrow \{b_2,b_1\}. \nonumber
\end{align}

\paragraph{Equation of motion relations}
The Bessel equation \eqref{eq:hdef} yields:
\begin{align}
	B[\{2,n_2,n_3,n_4\},\{a_1,a_2\},\{b_1,b_2\}] = & - B[\{0,n_2,n_3,n_4\},\{a_1,a_2\},\{b_1,b_2\}]\nonumber      \\
	                                               & - (2\nu+1) B[\{1,n_2,n_3,n_4\},\{a_1-1,a_2\},\{b_1+1,b_2\}], \\
	B[\{n_1,2,n_3,n_4\},\{a_1,a_2\},\{b_1,b_2\}] = & - B[\{n_1,0,n_3,n_4\},\{a_1,a_2\},\{b_1,b_2\}]\nonumber      \\
	                                               & - (2\nu+1) B[\{n_1,1,n_3,n_4\},\{a_1,a_2-1\},\{b_1+1,b_2\}], \\
	B[\{n_1,n_2,2,n_4\},\{a_1,a_2\},\{b_1,b_2\}] = & - B[\{n_1,n_2,0,n_4\},\{a_1,a_2\},\{b_1,b_2\}]\nonumber      \\
	                                               & - (2\nu+1) B[\{n_1,n_2,1,n_4\},\{a_1-1,a_2\},\{b_1,b_2+1\}], \\
	B[\{n_1,n_2,n_3,2\},\{a_1,a_2\},\{b_1,b_2\}] = & - B[\{n_1,n_2,n_3,0\},\{a_1,a_2\},\{b_1,b_2\}]\nonumber      \\
	                                               & - (2\nu+1) B[\{n_1,n_2,n_3,1\},\{a_1,a_2-1\},\{b_1,b_2+1\}],
\end{align}
where $\nu$ is the common mass parameter.

\paragraph{Symmetry relations}
1. \textbf{Propagator exchange}:
\begin{align}
	B[\{n_1,n_2,n_3,n_4\},\{a_1,a_2\},\{b_1,b_2\}] = B[\{n_3,n_4,n_1,n_2\},\{a_1,a_2\},\{b_2,b_1\}].
\end{align}

2. \textbf{Time exchange with energy swap}:
\begin{align}
	B^{(21)}[\{n_1,n_2,n_3,n_4\},\{a_1,a_2\},\{b_1,b_2\}; P_1,P_2] = B^{(12)}[\{n_2,n_1,n_4,n_3\},\{a_2,a_1\},\{b_1,b_2\}; P_2,P_1].
\end{align}
With the specialization $P_1=P_2$, this becomes a direct symmetry relation between the two time-orderings at fixed kinematics:
\begin{align}
	B[\{n_1,n_2,n_3,n_4\},\{a_1,a_2\},\{b_1,b_2\}; P_1,P_1] = B[\{n_2,n_1,n_4,n_3\},\{a_2,a_1\},\{b_1,b_2\}; P_1,P_1].
\end{align}

\paragraph{Differential operators in kinematics}
The dependence on external energies is generated by derivatives with respect to $P_1,P_2$:
\begin{align}
	\frac{\partial}{\partial P_1}B[\{n_1,n_2,n_3,n_4\},\{a_1,a_2\},\{b_1,b_2\};P_1,P_2] &= i\,B[\{n_1,n_2,n_3,n_4\},\{a_1+1,a_2\},\{b_1,b_2\};P_1,P_2], \label{eq:GdP1}\\
	\frac{\partial}{\partial P_2}B[\{n_1,n_2,n_3,n_4\},\{a_1,a_2\},\{b_1,b_2\};P_1,P_2] &= i\,B[\{n_1,n_2,n_3,n_4\},\{a_1,a_2+1\},\{b_1,b_2\};P_1,P_2]. \label{eq:GdP2}
\end{align}

\paragraph{Derivative with respect to $k_s$}
We use
\begin{align}
	\frac{\partial}{\partial k_s}|\bm q+\bm k_s|
	=\frac{k_s^2+\bm k_s\cdot \bm q}{k_s|\bm q+\bm k_s|}
	=\frac{k_s^2+|\bm q+\bm k_s|^2-q^2}{2k_s|\bm q+\bm k_s|},
\end{align}
which leads to a closed relation within the same $B$ family:
\begin{align}
	&\frac{\partial}{\partial k_s}B[\{n_1,n_2,n_3,n_4\},\{a_1,a_2\},\{b_1,b_2\}] \n\\
	&=-\frac{b_2}{2k_s}\Big(
	k_s^2\,B[\{n_1,n_2,n_3,n_4\},\{a_1,a_2\},\{b_1,b_2+2\}] \nonumber\\
	&\qquad
	+B[\{n_1,n_2,n_3,n_4\},\{a_1,a_2\},\{b_1,b_2\}] \nonumber\\
	&\qquad
	-B[\{n_1,n_2,n_3,n_4\},\{a_1,a_2\},\{b_1-2,b_2+2\}]
	\Big) \nonumber\\
	&+\frac{1}{2k_s}\Big(
	k_s^2\,B[\{n_1,n_2,n_3+1,n_4\},\{a_1+1,a_2\},\{b_1,b_2+1\}] \nonumber\\
	&\qquad
	+B[\{n_1,n_2,n_3+1,n_4\},\{a_1+1,a_2\},\{b_1,b_2-1\}] \nonumber\\
	&\qquad
	-B[\{n_1,n_2,n_3+1,n_4\},\{a_1+1,a_2\},\{b_1-2,b_2+1\}]
	\Big) \nonumber\\
	&+\frac{1}{2k_s}\Big(
	k_s^2\,B[\{n_1,n_2,n_3,n_4+1\},\{a_1,a_2+1\},\{b_1,b_2+1\}] \nonumber\\
	&\qquad
	+B[\{n_1,n_2,n_3,n_4+1\},\{a_1,a_2+1\},\{b_1,b_2-1\}] \nonumber\\
	&\qquad
	-B[\{n_1,n_2,n_3,n_4+1\},\{a_1,a_2+1\},\{b_1-2,b_2+1\}]
	\Big). \label{eq:Gdks}
\end{align}

\textbf{Scaling relation}:
\begin{align}
	\left( k_s \frac{\partial}{\partial k_s} + P_1 \frac{\partial}{\partial P_1} + P_2 \frac{\partial}{\partial P_2} \right) B[\{\bm n\},\{a_1,a_2\},\{b_1,b_2\}] \n\\
	= (D - b_1 - b_2 - a_1 - a_2 - 2) B[\{\bm n\},\{a_1,a_2\},\{b_1,b_2\}]. \label{eq:Gscaling}
\end{align}

\subsection{Tadpole families \texorpdfstring{$R_1,R_2$}{R1, R2}}
\paragraph{Tadpole families $R_1,R_2$}
The tadpole terms $R_1,R_2$ are single-time integrals generated by the time IBP coincidence terms. They require their own IBP system for reduction.

\paragraph{Kira check: whether $b_1$ can be reduced to $0$}
todo


\paragraph{IBP relations for $R_1$ (and $R_2$ via symmetry)}
The tadpole family $R_1$ satisfies the following IBP relations. The corresponding relations for $R_2$ follow via the symmetry $\bm q \to -(\bm q+\bm k_s)$, which exchanges $b_1\leftrightarrow b_2$ and relabels Hankel indices appropriately.

\textbf{Time IBP for $R_1$}:
\begin{align}
	0 = & -i(P_1+P_2) R_1[\{n_3,n_4\},\{a\},\{b_1,b_2\}] - a R_1[\{n_3,n_4\},\{a-1\},\{b_1,b_2\}] \nonumber       \\
	    & - R_1[\{n_3+1,n_4\},\{a\},\{b_1,b_2-1\}] - R_1[\{n_3,n_4+1\},\{a\},\{b_1,b_2-1\}]. \label{eq:R1timeIBP}
\end{align}

\textbf{Momentum IBP for $R_1$}:
from $\int \pd_{\bm q^i}[\bm q^i\cdots]=0$,
\begin{align}
	0 ={}& \int_{-\infty}^0 \di\tau\int \FR{\di^D \bm q}{(2\pi)^D}\,
	\Big[D+\mathcal{O}_1\Big]\Bigg[
	e^{-\ii P_0\tau}\frac{(-\tau)^{a}}{x_1^{b_1}x_2^{b_2}}\,
	h_{n_3}^{(1)}(-x_2\tau)h_{n_4}^{(2)}(-x_2\tau)\Bigg] \nonumber\\
	={}& \int_{-\infty}^0 \di\tau\int \FR{\di^D \bm q}{(2\pi)^D}\,
	\Big[D+(\mathcal{O}_1 x_1)\partial_{x_1}+(\mathcal{O}_1 x_2)\partial_{x_2}\Big]\Bigg[
	e^{-\ii P_0\tau}\frac{(-\tau)^{a}}{x_1^{b_1}x_2^{b_2}}\,
	h_{n_3}^{(1)}(-x_2\tau)h_{n_4}^{(2)}(-x_2\tau)\Bigg]. \label{eq:R1momIBP_chainrule}
\end{align}
\begin{align}
	0 = & D R_1[\{n_3,n_4\},\{a\},\{b_1,b_2\}] - b_1 R_1[\{n_3,n_4\},\{a\},\{b_1,b_2\}] \nonumber                    \\
	    & - \FR12 b_2\Big( R_1[\{n_3,n_4\},\{a\},\{b_1,b_2\}] + R_1[\{n_3,n_4\},\{a\},\{b_1-2,b_2+2\}] \nonumber     \\
	    & -k_s^2 R_1[\{n_3,n_4\},\{a\},\{b_1,b_2+2\}] \Big) \nonumber                                                \\
	    & + \FR12 \Big(
	    R_1[\{n_3+1,n_4\},\{a+1\},\{b_1,b_2-1\}] + R_1[\{n_3+1,n_4\},\{a+1\},\{b_1-2,b_2+1\}] \nonumber\\
	    &\qquad\qquad\qquad\quad -k_s^2 R_1[\{n_3+1,n_4\},\{a+1\},\{b_1,b_2+1\}]
	    \Big) \nonumber\\
	    & + \FR12 \Big(
	    R_1[\{n_3,n_4+1\},\{a+1\},\{b_1,b_2-1\}] + R_1[\{n_3,n_4+1\},\{a+1\},\{b_1-2,b_2+1\}] \nonumber\\
	    &\qquad\qquad\qquad\quad -k_s^2 R_1[\{n_3,n_4+1\},\{a+1\},\{b_1,b_2+1\}]
	    \Big). \label{eq:R1momIBP}
\end{align}

\textbf{Complementary momentum IBP for $R_1$}:
from $\int \pd_{\bm q^i}[(\bm q^i+\bm k_s^i)\cdots]=0$,
\begin{align}
	0 ={}& \int_{-\infty}^0 \di\tau\int \FR{\di^D \bm q}{(2\pi)^D}\,
	\Big[D+\mathcal{O}_2\Big]\Bigg[
	e^{-\ii P_0\tau}\frac{(-\tau)^{a}}{x_1^{b_1}x_2^{b_2}}\,
	h_{n_3}^{(1)}(-x_2\tau)h_{n_4}^{(2)}(-x_2\tau)\Bigg] \nonumber\\
	={}& \int_{-\infty}^0 \di\tau\int \FR{\di^D \bm q}{(2\pi)^D}\,
	\Big[D+(\mathcal{O}_2 x_1)\partial_{x_1}+(\mathcal{O}_2 x_2)\partial_{x_2}\Big]\Bigg[
	e^{-\ii P_0\tau}\frac{(-\tau)^{a}}{x_1^{b_1}x_2^{b_2}}\,
	h_{n_3}^{(1)}(-x_2\tau)h_{n_4}^{(2)}(-x_2\tau)\Bigg]. \label{eq:R1momIBP2_chainrule}
\end{align}
\begin{align}
	0 = & D R_1[\{n_3,n_4\},\{a\},\{b_1,b_2\}] - b_2 R_1[\{n_3,n_4\},\{a\},\{b_1,b_2\}] \nonumber                    \\
	    & - \FR12 b_1\Big( R_1[\{n_3,n_4\},\{a\},\{b_1,b_2\}] + R_1[\{n_3,n_4\},\{a\},\{b_1+2,b_2-2\}] \nonumber     \\
	    & -k_s^2 R_1[\{n_3,n_4\},\{a\},\{b_1+2,b_2\}] \Big) \nonumber                                                \\
	    & + R_1[\{n_3+1,n_4\},\{a+1\},\{b_1,b_2-1\}] + R_1[\{n_3,n_4+1\},\{a+1\},\{b_1,b_2-1\}]. \label{eq:R1momIBP2}
\end{align}

\textbf{Equation of motion for $R_1$}:
\begin{align}
	R_1[\{2,n_4\},\{a\},\{b_1,b_2\}] = & - R_1[\{0,n_4\},\{a\},\{b_1,b_2\}] \nonumber                            \\
	                                   & - (2\nu+1) R_1[\{1,n_4\},\{a-1\},\{b_1,b_2+1\}], \label{eq:R1eom1} \\
	R_1[\{n_3,2\},\{a\},\{b_1,b_2\}] = & - R_1[\{n_3,0\},\{a\},\{b_1,b_2\}] \nonumber                            \\
	                                   & - (2\nu+1) R_1[\{n_3,1\},\{a-1\},\{b_1,b_2+1\}]. \label{eq:R1eom2}
\end{align}

\textbf{Symmetry relations}: $R_1$ inherits symmetry from the bubble:
\begin{align}
	 R_1[\{n_3,n_4\},\{a\},\{b_1,b_2\}] = R_1[\{n_4,n_3\},\{a\},\{b_1,b_2\}] \label{eq:R1sym}. 
\end{align}	 
	 More importantly, $R_1$ and $R_2$ are related by loop-momentum shift $\bm q \to -(\bm q+\bm k_s)$:
\begin{align}
	R_2[\{n_1,n_2\},\{a\},\{b_1,b_2\}] = R_1[\{n_1,n_2\},\{a\},\{b_2,b_1\}]. \label{eq:R1R2sym}
\end{align}
This symmetry implies that one only needs to implement the IBP system for $R_1$; all relations for $R_2$ follow by applying $b_1\leftrightarrow b_2$ and appropriate index relabeling.

These tadpole integrals are simpler than the double-time $B$ integrals because they involve only one time variable $\tau$ and two Hankel functions (vs four in $B$), and can be reduced to master integrals using the above IBP system.



\paragraph{Differential operators for $R_1$}
From the definition of $R_1$, we can derive relations by taking derivatives with respect to parameters:

\textbf{Derivative with respect to $P_0$}:
\begin{align}
	\frac{\partial}{\partial P_0} R_1[\{n_3,n_4\},\{a\},\{b_1,b_2\}] = i R_1[\{n_3,n_4\},\{a+1\},\{b_1,b_2\}]. \label{eq:R1dP0}
\end{align}

\textbf{Derivative with respect to $k_s$}:
\begin{align}
	&\frac{\partial}{\partial k_s} R_1[\{n_3,n_4\},\{a\},\{b_1,b_2\}] \n\\
	&=-\frac{b_2}{2k_s}\Big(
	k_s^2\,R_1[\{n_3,n_4\},\{a\},\{b_1,b_2+2\}] \nonumber\\
	&\qquad
	+R_1[\{n_3,n_4\},\{a\},\{b_1,b_2\}] \nonumber\\
	&\qquad
	-R_1[\{n_3,n_4\},\{a\},\{b_1-2,b_2+2\}]
	\Big) \nonumber\\
	&+\frac{1}{2k_s}\Big(
	k_s^2\,R_1[\{n_3+1,n_4\},\{a+1\},\{b_1,b_2+1\}] \nonumber\\
	&\qquad
	+R_1[\{n_3+1,n_4\},\{a+1\},\{b_1,b_2-1\}] \nonumber\\
	&\qquad
	-R_1[\{n_3+1,n_4\},\{a+1\},\{b_1-2,b_2+1\}]
	\Big) \nonumber\\
	&+\frac{1}{2k_s}\Big(
	k_s^2\,R_1[\{n_3,n_4+1\},\{a+1\},\{b_1,b_2+1\}] \nonumber\\
	&\qquad
	+R_1[\{n_3,n_4+1\},\{a+1\},\{b_1,b_2-1\}] \nonumber\\
	&\qquad
	-R_1[\{n_3,n_4+1\},\{a+1\},\{b_1-2,b_2+1\}]
	\Big). \label{eq:R1dks}
\end{align}

\textbf{Scaling relation} (dimensional analysis):
\begin{align}
	\left( k_s \frac{\partial}{\partial k_s} + P_0 \frac{\partial}{\partial P_0} \right) R_1[\{n_3,n_4\},\{a\},\{b_1,b_2\}] = (D - b_1 - b_2 - a - 1) R_1[\{n_3,n_4\},\{a\},\{b_1,b_2\}]. \label{eq:R1scaling}
\end{align}
This scaling relation will be used to verify the correctness of differential equations.





\subsection{Odevity seclection rules for the IBP system}
There are two useful parity classifications (``Odevity'') for the bubble IBP system:
\begin{enumerate}
\item For $\di q$-IBP and $\di \tau$-IBP, the parities of $n_1+n_2+b_1$ and $n_3+n_4+b_2$ are conserved.
\item For $\di q$-IBP, the parities of $n_1+n_3+a_1$ and $n_2+n_4+a_2$ are conserved.
\end{enumerate}
In reduction runs we restrict the IBP system to the even sector of the first classification, i.e.
\begin{align}
	n_1+n_2+b_1 \ \text{even},\qquad n_3+n_4+b_2 \ \text{even}. \label{eq:odevity_even_sector}
\end{align}
For the tadpole families this implies two independent constraints:
\begin{itemize}
\item For $R_1$ (obtained by contracting indices where one is 1 and the other is 0), one must impose $b_1$ even \emph{and} $n_3+n_4+b_2$ even. This contraction produces a term with an odd power (excluding the cancelling $2\nu$ shift), which matches the Odevity of the top sector.
\item For $R_2$, the loop-momentum shift symmetry \eqref{eq:R1R2sym} exchanges $b_1\leftrightarrow b_2$, so the corresponding restriction is $b_2$ even and $n_1+n_2+b_1$ even.
\end{itemize}
This is why Odevity belongs to the IBP chapter: it specifies which integrals are allowed to appear in the closed IBP system that we export to Kira.


\subsection{Reduction strategy}
\begin{enumerate}
	\item Select a Odevity sub-system of IBP-system. Apply EoM and symmetry to time-IBP and momentum-IBP, Then feed the processed IBP relations to Kira for reduction, since Mathematica is too slow for the reduction step.
	\item Reduce $R_1,R_2$ first. then, reduce $B$ with the known $R_1,R_2$ in time IBP.
\end{enumerate}


\section{Baikov representation and dimension-shift relations}
\label{sec:baikov}
\paragraph{Baikov variables for the one-loop two-point momentum integral}
Introduce $x_1\equiv|\bm q|$, $x_2\equiv|\bm q+\bm k_s|$ and
\begin{align}
	z_1=x_1^2=q^2,\qquad z_2=x_2^2=(q+k_s)^2,\qquad \mK=k_s^2,
\end{align}
with Gram determinant
\begin{align}
	\mG=\frac{1}{4}\Big(2 z_1 k_s^2+2 z_2 k_s^2-k_s^4-z_1^2-z_2^2+2 z_1 z_2\Big)
	=\frac{1}{4}\lambda(z_1,z_2,k_s^2),
\end{align}
where $\lambda$ is the K\"all\'en function. The one-loop Baikov change of variables gives (for general $d$)
\begin{align}
	\di^d q \to C_d\,\mG^{\frac{d-3}{2}}\mK^{-\frac{d-2}{2}}\,\di z_1\,\di z_2,
	\qquad C_{d} = \frac{\pi^{-1/2}}{\Gamma[(d-1)/2]}.
\end{align}
In particular, since $\di z_i=2x_i\,\di x_i$, one has the useful bookkeeping relation
\begin{align}
	B[\{\bm{n}\},\{a_1,a_2\},\{b_1+1,b_2+1\}]
	\sim \cdots \times \frac{\di z_1}{x_1^{b_0+b_1+1}}\frac{\di z_2}{x_2^{b_0+b_2+1}}
	= \cdots \times 4\,\frac{\di x_1}{x_1^{b_0+b_1}}\frac{\di x_2}{x_2^{b_0+b_2}}.
\end{align}
The integration region is determined by $\mG\ge 0$, equivalently the triangle inequalities between $(x_1,x_2,k_s)$.

\paragraph{Gram matrices and dimension-shift identity}
For a general $L$-loop integral with $E$ independent external momenta, define the Gram matrices
\begin{align}
	(\mG_{\text{full}})_{ij} \equiv q_i \cdotp q_j, \qquad \{q_i\} = \{l_1,\dots,l_L,p_1,\dots,p_E\},
\end{align}
where $l_i$ are loop momenta and $p_i$ are external momenta. The Gram determinant is $\mG \equiv \det(\mG_{\text{full}})$. The kinematic Gram matrix involves only external momenta:
\begin{align}
	\mK_{ij} \equiv p_i \cdotp p_j, \qquad i,j = 1,\dots,E,
\end{align}
with $\mK \equiv \det(\mK_{ij})$.

\paragraph{Dimension shift}
With generalized Sylvester's determinant identity, one has
\begin{align}
	|\partial_{(q_i\cdotp q_j)}|_{i,j\leq L} |q_i\cdotp q_j|^\gamma
	= |q_i\cdotp q_j|_{i,j> L}\left(\prod_{i=0}^{L-1} (\gamma+i)\right)|q_i\cdotp q_j|^{\gamma-1}.
\end{align}
For a symmetric matrix of scalar products one further has
\begin{align}
	&|A_{ij} \partial_{(q_i\cdotp q_j)}|_{i,j\leq L} |q_i\cdotp q_j|^\gamma = |q_i\cdotp q_j|_{i,j> L}  \left(\prod_{i=0}^{L-1} \left(\gamma+\frac{i}{2}\right)\right)   |q_i\cdotp q_j|^{\gamma-1}  \n\\
	&A_{ij}=\frac{1}{2} ~~\text{for}~~i\neq j,  ~~A_{ij}=1 ~~\text{for}~~i=j.
\end{align}
Here $|q_i\cdotp q_j|_{i,j> L} = \mK$ is the kinematic Gram determinant, and $|q_i\cdotp q_j| = \mG$ is the full Gram determinant. Writing
\begin{align}
	Q=\frac{1}{\prod_i z_i^{a_i}},\qquad \gamma=\frac{d-E-L-1}{2},
\end{align}
the identity above leads to a dimension-raising relation in Baikov form:
\begin{align}
	C_{\gamma-1}\!\int Q\,\frac{\mG^{\gamma-1}}{\mK^{\gamma-1+L/2}}\prod \di z_i
	=&~C_{\gamma}\!\int Q\,|A_{ij} \partial_{(q_i\cdotp q_j)}|_{i,j\leq L}\,
	\frac{\mG^\gamma}{\mK^{\gamma+L/2}}\prod \di z_i \n\\
	=&~(-1)^L C_{\gamma}\!\int Q\,|\overleftarrow{A_{ij} \partial_{(q_i\cdotp q_j)}}|_{i,j\leq L}\,
	\frac{\mG^\gamma}{\mK^{\gamma+L/2}}\prod \di z_i,
\end{align}
where the left-hand side corresponds to a $d-2$ dimensional integral while the right-hand side produces a $d$ dimensional one after integrating by parts.

\section{Intersection theory viewpoint and dlog integrands}
\label{sec:dlog}

\paragraph{Convention}
In this section, we use the notation $\hat{F}$ to denote the integrand of a function $F$. For example, $\hat{B}$ is the integrand of the double-time integral $B$, and $\hat{I}$ is the integrand of the time integral $I$.

We define the shift parameters $a_0$ and $b_0$ based on the regularization and Hankel-function asymptotics.
For the top-sector Bubble $B$, the time variables $\tau_1, \tau_2$ each carry a regularization index $\delta_a$. Additionally, we extract a factor $(-x_i \tau_i)^{-\nu}$ from each Hankel function to isolate the non-integer powers. With two Hankel functions per time variable, this adds $2\nu$ to the $\tau$ weight and $-2\nu$ to the spatial weight. Thus we define:
\begin{align}
	a_0 &= \delta_a + 2\nu \quad (\text{for } B), \n\\
	b_0 &= -2\nu.
\end{align}
For the Tadpole family $R_1$, the structure arises from contracting one propagator in the Bubble (effectively setting $\tau_1=\tau_2=\tau$). The regularization indices sum up, $\delta_a+\delta_a=2\delta_a$. The surviving propagator contributes the standard Hankel shift of $2\nu$. Thus for $R_1$ we use:
\begin{align}
	a_0 &= 2\delta_a + 2\nu \quad (\text{for } R_1), \n\\
	b_0 &= -2\nu.
\end{align}
In the following, the symbol $a_0$ refers to the shift parameter appropriate for the specific integral family under discussion.

By intersection theory, if a master-integrand basis $\phi_i$ is of $d\log$ form and the twist $u$ satisfies $\di u=(\di\log u)\,\nu$, then the differential equations for the master integrals $\int \phi_i u$ take a $d\log$ form. Here we use a mild generalization motivated by fibration intersection theory: treat the master objects as $\Phi_i\cdot U$ with both $\Phi_i$ and $U$ viewed as vectors, choose $U$ to satisfy
\begin{align}
	\di U = \di\Omega \cdot U,
\end{align}
and construct $\Phi_i$ as $d\log$-form integrands. This motivates a canonical-like choice of masters for the de Sitter bubble.

\paragraph{Time-integral block as a first-order system}
Define the ``tree-level-like'' time integral (no momentum denominators)
\begin{align}
	I_{\bm n}(x_1,x_2;P_1,P_2)\equiv \int_{-\infty}^0 \di\tau_1\di\tau_2\,\hat{I}_{\bm n},
\end{align}
with the integrand
\begin{align}
	\hat{I}_{\bm n} \equiv & e^{-\ii P_1\tau_1-\ii P_2\tau_2}\,(-\tau_1)^{a_0}(-\tau_2)^{a_0} \nonumber\\
	&\times h_{n_1}^{(1)}(-x_1\tau_1)h_{n_2}^{(2)}(-x_1\tau_2)\,h_{n_3}^{(1)}(-x_2\tau_1)h_{n_4}^{(2)}(-x_2\tau_2).
\end{align}
For fixed $(\nu,a_0)$, the $I_{\bm n}$ span a finite basis and satisfy a first-order system
\begin{align}
	\di I_{\bm m}=(\di\Omega_{\bm m\bm n})\,I_{\bm n},
\end{align}
where the one-forms in $\di\Omega$ are built from $\di\log x_1$, $\di\log x_2$ and $\di\log(P_i\pm x_1\pm x_2)$.
For the two-vertex tree-level-like family entering the bubble, there are $2^4$ time-integral master functions.
\begin{align}
	\di\Omega=\di\Omega_{x_1}+\di\Omega_{x_2}+\di\Omega_{ex},
\end{align}
where $\di\Omega_{x_i}$ only contains $\di\log x_i$ while $\di\Omega_{ex}$ is composed of $\di\log(P_i\pm x_1\pm x_2)$ with constant coefficients (built from $\nu$ and $a_0$).

\paragraph{Baikov twist and dlog building blocks}
With $\ep$ the dimensional regulator ($d=3-2\ep$), the one-loop Baikov measure contributes the factor $\mG^{-\ep}\mK^{\ep}$, and one can package the combined twist schematically as a vector
\begin{align}
	U = x_1^{-b_0}x_2^{-b_0}\,\mG^{-\ep}\mK^{\ep}\,I_{\bm n},
\end{align}
whose differential is again a sum of $d\log$ one-forms,
\begin{align}
	\di U=\Big(\di\Omega-b_0(\di\log x_1+\di\log x_2)-\ep(\di\log \mG-\di\log \mK)\Big)\cdot U.
\end{align}
In terms of $z_i=x_i^2$, useful $d\log$ building blocks include
\begin{align}
	&\di\log z_1=2\di\log x_1,\qquad \di\log z_2=2\di\log x_2,\n\\
	&\frac{\sqrt{z_2 k_s^2}\,\di z_1}{\mG}
	= - \di \log \left(\frac{2 \sqrt{z_2 k_s^2}+k_s^2-z_1+z_2}{-2 \sqrt{z_2 k_s^2}+k_s^2-z_1+z_2}\right),
	\qquad z_1\leftrightarrow z_2.
\end{align}
Since $\di\!\int \di x_i/(P_i\pm x_1\pm x_2)$ is of $d\log$ form, one also has the building blocks
\begin{align}
	&\partial_{P_j}\Omega_{ex}\,\di x_i = \partial_{P_j}\Omega\,\di x_i, \n\\
	&\partial_{x_j}\Omega_{ex}\,\di x_i = (\partial_{x_j}\Omega-\partial_{x_j}\Omega_{x_j})\,\di x_i.
\end{align}
Using
\begin{align}
	\hat{B}[\{\bm{n}\},\{a_1+1,a_2\},\{b_1,b_2\}]
	&=\partial_{P_1}\hat{B}[\{\bm{n}\},\{a_1,a_2\},\{b_1,b_2\}]
	=\left(\partial_{P_1}\Omega_{\bm n\bm m}\right)\,\n\\
	&\times \hat{B}[\{\bm{m}\},\{a_1,a_2\},\{b_1,b_2\}], \n\\
	\hat{B}[\{\bm{n}\},\{a_1,a_2\},\{b_1,b_2\}]
	&=\frac{1}{b_0+b_1-1}\left(\partial_{x_1}\Omega_{\bm n\bm m}\right)\,\n\\
	&\times \hat{B}[\{\bm{m}\},\{a_1,a_2\},\{b_1-1,b_2\}],
\end{align}
one may view $\partial_{P_j}\Omega\,\di x_i\sim \tau_j\,\di x_i$ and
$\partial_{x_j}\Omega\,\di x_i\sim \frac{b_0+b_j-1}{x_j}\di x_j$
as additional $d\log$-type building blocks at the integrand level.

\subsection{Top-sector \texorpdfstring{$d\log$}{dlog} candidates}
\paragraph{Constructing candidate $d\log$ integrands}
For notational convenience, we split off common shifts and write
\begin{align}
	\hat{B}[\{\bm n\},\{a_0+a_1,a_0+a_2\},\{b_0+b_1,b_0+b_2\}]\to \hat{B}[\{\bm n\},\{a_1,a_2\},\{b_1,b_2\}],
\end{align}
where $a_i,b_i\in\mathbb Z$ and $(a_0,b_0)$ encode the universal shifts induced by the special-function normalization and prefactors.

Some candidate $d\log$-type integrands in $d=3-2\ep$ are (up to overall constants)
\begin{align}
	k_s\,\hat{B}[\{\bm n\},\{0,0\},\{2,2\};d=3]
	&= \frac{\di x_1\,\di x_2}{x_1 x_2}\,x_1^{-b_0}x_2^{-b_0}\,I^{(R)}_{\bm n}\,\mG^{-\ep}\mK^{\ep}, \n\\
	k_s\,\hat{B}[\{\bm n\},\{1,0\},\{2,1\};d=3]
	&= \frac{\di x_1}{x_1}\,\tau_1\,\di x_2\,x_1^{-b_0}x_2^{-b_0}\,I^{(R)}_{\bm n}\,\mG^{-\ep}\mK^{\ep}, \n\\
	k_s\,\hat{B}[\{\bm n\},\{1,1\},\{1,1\};d=3]
	&= \tau_1\,\di x_1\,\tau_2\,\di x_2\,x_1^{-b_0}x_2^{-b_0}\,I^{(R)}_{\bm n}\,\mG^{-\ep}\mK^{\ep}\qquad \text{\cjq{wrong, double pole}}.
\end{align}
In $d=1-2\ep$, one can also use the Baikov building block $\frac{k_s\sqrt{z_1}\,\di z_2}{\mG}$ to form $d\log$ candidates such as
\begin{align}
	\hat{B}[\{\bm n\},\{0,0\},\{1,0\};d=1]
	&= \frac{\di z_1}{z_1}\,\frac{k_s\sqrt{z_1}\,\di z_2}{\mG}\,
	z_1^{-b_0/2}z_2^{-b_0/2}\,I^{(R)}_{\bm n}\,\mG^{-\ep}\mK^{\ep}, \n\\
	\hat{B}[\{\bm n\},\{1,0\},\{0,0\};d=1]
	&= \tau_1\,\di x_1\,\frac{k_s\sqrt{z_1}\,\di z_2}{\mG}\,
	z_1^{-b_0/2}z_2^{-b_0/2}\,I^{(R)}_{\bm n}\,\mG^{-\ep}\mK^{\ep}.
\end{align}
In addition, one may also consider candidates generated by $\Omega_{ex}$:
\begin{align}
	(\partial_{x_2}\Omega_{ex})_{\bm n \bm m}\,\di x_1\,\frac{\di x_2}{x_2}\,
	x_1^{-b_0}x_2^{-b_0}\,I^{(R)}_{\bm m}\,\mG^{-\ep}\mK^{\ep}, \n\\
	(\partial_{x_2}\Omega_{ex})_{\bm n \bm m}\,\di x_1\,\frac{k_s\sqrt{z_1}\,\di z_2}{\mG}\,
	z_1^{-b_0/2}z_2^{-b_0/2}\,I^{(R)}_{\bm m}\,\mG^{-\ep}\mK^{\ep}.
\end{align}

\paragraph{Master-integral selection for the bubble}
Based on the $d\log$ construction and IBP reduction, we select the following independent master integrals for the one-loop two-point bubble. The selection is performed at the original level (1-dimensional, without dimension raising), yielding 14 master integrals organized by dimension and index structure.

\textbf{Master integrals at $d=1$} (5 integrals):
\begin{align}
	& B[\{0, 0, 0, 1\}, \{0, 0\}, \{0, 1\}, 1], \quad
	B[\{1, 1, 0, 1\}, \{0, 0\}, \{0, 1\}, 1], \nonumber\\
	& B[\{0, 0, 0, 0\}, \{0, 1\}, \{0, 0\}, 1], \quad
	B[\{0, 0, 1, 1\}, \{0, 1\}, \{0, 0\}, 1], \nonumber\\
	& B[\{1, 1, 1, 1\}, \{0, 1\}, \{0, 0\}, 1].
\end{align}

\textbf{Master integrals at $d=3$} (9 integrals):

Type A --$\{0,1\}, \{1,2\}$ index pattern:
\begin{align}
	& B[\{0, 1, 0, 0\}, \{0, 1\}, \{1, 2\}, 3], \quad
	B[\{0, 1, 1, 1\}, \{0, 1\}, \{1, 2\}, 3], \nonumber\\
	& B[\{1, 0, 0, 0\}, \{0, 1\}, \{1, 2\}, 3], \quad
	B[\{1, 0, 1, 1\}, \{0, 1\}, \{1, 2\}, 3].
\end{align}

Type B --$\{0,0\}, \{2,2\}$ $d\log$ form:
\begin{align}
	& B[\{0, 0, 0, 0\}, \{0, 0\}, \{2, 2\}, 3], \quad
	B[\{0, 0, 1, 1\}, \{0, 0\}, \{2, 2\}, 3], \quad
	B[\{1, 1, 1, 1\}, \{0, 0\}, \{2, 2\}, 3].
\end{align}
These three integrals have the natural $d\log$ structure $k_s\,\hat{B}[\{n\},\{0,0\},\{2,2\};d=3] = \frac{\di x_1\,\di x_2}{x_1 x_2} \cdots$.

Type C --linear combinations with prefactor $(2+b_0+2\nu)$:
\begin{align}
	& (2 + b_0 + 2\nu)\,B[\{0, 1, 0, 1\}, \{0, 0\}, \{1, 3\}, 3], \nonumber\\
	& (2 + b_0 + 2\nu)\,B[\{0, 1, 1, 0\}, \{0, 0\}, \{1, 3\}, 3].
\end{align}
The prefactor $(2+b_0+2\nu)$ arises from the coefficient normalization during IBP reduction.

This master-integral basis is recorded in \texttt{codebubble/kira\_1L2p/result/MIdlog.m} and serves as the starting point for constructing differential equations in external kinematics.

\paragraph{Raising $d=1$ back to $d=3$}
For the Baikov measure one has
\begin{align}
	\di^d q \to C_d\,\mG^{\frac{d-3}{2}}\mK^{-\frac{d-2}{2}}\,\di z_1\,\di z_2.
\end{align}
In particular, for $d=1-2\ep$ and $d=3-2\ep$,
\begin{align}
	C_{1-2\ep}\,\mG^{-\ep-1}\mK^{\ep+1/2}
	= \partial_{q^2}\!\left[C_{3-2\ep}\,\mK^{-1/2+\ep}\mG^{-\ep}\right],
\end{align}
where $\partial_{q^2}\equiv\partial_{z_1}+\partial_{z_2}$ and we used $\partial_{q^2}\mG=\mK$ together with $\Gamma(1-\ep)=-\ep\,\Gamma(-\ep)$.

Therefore, for
\begin{align}
	\hat{B}[\{\bm n\},\{a_1,a_2\},\{b_1,b_2\};d=1]
	=&~\di z_1\,\di z_2\,z_1^{-(b_0+b_1)/2}z_2^{-(b_0+b_2)/2}\,I^{(R)}_{\bm n}\,\n\\
	&\times (-\tau_1)^{a_1}(-\tau_2)^{a_2}\,C_{1-2\ep}\,\mG^{-\ep-1}\mK^{\ep+1/2},
\end{align}
integration by parts gives
\begin{align}
	\hat{B}[\{\bm n\},\{a_1,a_2\},\{b_1,b_2\};d=1]
	=&~\di z_1\,\di z_2\,\Bigg[\left(\frac{b_0+b_1}{2z_1}+\frac{b_0+b_2}{2z_2}\right)I^{(R)}_{\bm n}
	-\left(\partial_{z_1}+\partial_{z_2}\right)I^{(R)}_{\bm n}\Bigg] \n\\
	&\times (-\tau_1)^{a_1}(-\tau_2)^{a_2}\,z_1^{-(b_0+b_1)/2}z_2^{-(b_0+b_2)/2}\,
	C_{3-2\ep}\,\mG^{-\ep}\mK^{-1/2+\ep}.
\end{align}
Equivalently, in $(x_1,x_2)$ variables with $\partial_{q^2}=\frac{1}{2x_1}\partial_{x_1}+\frac{1}{2x_2}\partial_{x_2}$,
\begin{align}
	\hat{B}[\{\bm n\},\{a_1,a_2\},\{b_1,b_2\};d=1]
	=&~2\,\di x_1\,\di x_2\,\Bigg[\left(\frac{b_0+b_1}{x_1^2}+\frac{b_0+b_2}{x_2^2}\right)I^{(R)}_{\bm n}
	-\left(\frac{1}{x_1}\partial_{x_1}+\frac{1}{x_2}\partial_{x_2}\right)I^{(R)}_{\bm n}\Bigg] \n\\
	&\times (-\tau_1)^{a_1}(-\tau_2)^{a_2}\,x_1^{-(b_0+b_1-1)}x_2^{-(b_0+b_2-1)}\,
	C_{3-2\ep}\,\mG^{-\ep}\mK^{-1/2+\ep}.
\end{align}
In particular, for the $d=1$ candidates above,
\begin{align}
	\hat{B}[\{\bm n\},\{0,0\},\{1,0\};d=1]
	&= \di z_1\,\di z_2\,z_1^{-b_0/2-1/2}z_2^{-b_0/2}\,I^{(R)}_{\bm n}\,\mG^{-\ep-1}\mK^{\ep+1/2} \n\\
	&=2\,\di x_1\,\di x_2\,\Bigg[\left(\frac{b_{0}+1}{x_1^2}+\frac{b_{0}}{x_2^2}\right)I^{(R)}_{\bm n}
	-\left(\frac{1}{x_1}\partial_{x_1}+\frac{1}{x_2}\partial_{x_2}\right)I^{(R)}_{\bm n}\Bigg]\n\\
	&\times x_1^{-b_0}x_2^{-b_0+1}\,\mG^{-\ep}\mK^{-1/2+\ep}, \n\\
	\hat{B}[\{\bm n\},\{1,0\},\{0,0\};d=1]
	&=2\,\di x_1\,\di x_2\,\Bigg[b_0\left(\frac{1}{x_1^2}+\frac{1}{x_2^2}\right)I^{(R)}_{\bm n}
	-\left(\frac{1}{x_1}\partial_{x_1}+\frac{1}{x_2}\partial_{x_2}\right)I^{(R)}_{\bm n}\Bigg]\n\\
	&\times (-\tau_1)\,x_1^{-b_0+1}x_2^{-b_0+1}\,\mG^{-\ep}\mK^{-1/2+\ep}.
\end{align}
Moreover,
\begin{align*}
	&(\partial_{x_2}\Omega_{ex})_{\bm n \bm m}\,\di x_1\,\frac{\di x_2}{x_2}\,x_1^{-b_0}x_2^{-b_0}\,I^{(R)}_{\bm m}\,\mG^{-\ep}\mK^{\ep} \n\\
	&= (\partial_{x_2}\Omega-\partial_{x_2}\Omega_{x_2})_{\bm n \bm m}\,\di x_1\,\frac{\di x_2}{x_2}\,x_1^{-b_0}x_2^{-b_0}\,I^{(R)}_{\bm m}\,\mG^{-\ep}\mK^{\ep} \n\\
	&= \di x_1\,\di x_2\,(\partial_{x_2}I^{(R)}_{\bm n})\,x_1^{-b_0}x_2^{-b_0-1}\,\mG^{-\ep}\mK^{\ep} 
	- \di x_1\,\di x_2\,(\partial_{x_2}\Omega_{x_2})_{\bm n \bm m}\,x_1^{-b_0}x_2^{-b_0-1}\,I^{(R)}_{\bm m}\,\mG^{-\ep}\mK^{\ep} \n\\
	&= \di x_1\,(\di I^{(R)}_{\bm n})\,x_1^{-b_0}x_2^{-b_0-1}\,\mG^{-\ep}\mK^{\ep} 
	- \di x_1\,\di x_2\,(\partial_{x_2}\Omega_{x_2})_{\bm n \bm m}\,x_1^{-b_0}x_2^{-b_0-1}\,I^{(R)}_{\bm m}\,\mG^{-\ep}\mK^{\ep} \n\\
	&= -\di x_1\,I^{(R)}_{\bm n}\,\di\!\left( x_1^{-b_0}x_2^{-b_0-1}\,\mG^{-\ep}\mK^{\ep} \right) 
	- \di x_1\,\di x_2\,(\partial_{x_2}\Omega_{x_2})_{\bm n \bm m}\,x_1^{-b_0}x_2^{-b_0-1}\,I^{(R)}_{\bm m}\,\mG^{-\ep}\mK^{\ep} \n\\
	&= \left(\frac{b_0+1}{x_2}\delta_{\bm n \bm m}-(\partial_{x_2}\Omega_{x_2})_{\bm n \bm m}\right)\,k_s\,x_1^{-b_0}x_2^{-b_0-1}\,I^{(R)}_{\bm m}\,\mG^{-\ep}\mK^{\ep-1/2}\,\di x_1\,\di x_2.
\end{align*}
and similarly
\begin{align*}
	&(\partial_{x_2}\Omega_{ex})_{\bm n \bm m}\,\di x_1\,\frac{k_s\sqrt{z_1}\,\di z_2}{\mG}\,z_1^{-b_0/2}z_2^{-b_0/2}\,I^{(R)}_{\bm m}\,\mG^{-\ep}\mK^{\ep} \n\\
	&= \left(\frac{b_0-1}{x_2}\delta_{\bm n \bm m}-(\partial_{x_2}\Omega_{x_2})_{\bm n \bm m}\right)\,x_1^{-b_0+1}x_2^{-b_0+1}\,I^{(R)}_{\bm m}\,\mG^{-\ep-1}\mK^{\ep+1/2}\,\di x_1\,\di x_2.
\end{align*}

The practical goal is to pick a master basis for the full de Sitter bubble (including the special-function time dependence) such that the induced differential system in external kinematics is as close as possible to a $d\log$ form.

\subsection{Alphabet and differential system structure}
Based on the intersection-theory analysis, we have computed the differential equations for the master integrals.
In the $d\log$ form, the differential system is characterized by an alphabet of letters.
Unlike in flat-space calculations where the alphabet typically depends only on kinematic variables, in the massive de Sitter case with V++ vertices, the letters exhibit a non-trivial dependence on the mass parameter $\nu$ and the dimensional regulator $\epsilon$.

The alphabet $\mathcal{A}$ appearing in the differential equations with respect to the kinematic variable $x$ is found to be:
\begin{align}
	\mathcal{A} = \left\{ x, x\pm 1, x\pm \sqrt{1+2\epsilon}, x \pm \frac{1+2\nu \pm \sqrt{3+4\epsilon+4\nu(1+\nu)}}{\sqrt{2}} \right\}.
\end{align}
We observe that the alphabet contains square roots involving $\epsilon$ and $\nu$, such as $\sqrt{1+2\epsilon}$ and $\sqrt{3+4\epsilon+4\nu(1+\nu)}$.
This deviation from the rational structure usually found in flat-space polylogarithms is a hallmark of the massive de Sitter propagators and the specific time-ordering structure.

The coefficient matrix of the differential system, while too lengthy to be presented explicitly here, is constructed from rational functions of $\epsilon$ and $\nu$, but decorated with the aforementioned square-root structures.
These results for the alphabet and the structure of the differential equations provide the necessary data for expressing the master integrals in terms of iterated integrals, potentially generalizing the standard multiple polylogarithms to include algebraic kernels.





\subsection{Tadpole-family \texorpdfstring{$d\log$}{dlog} candidates}
\label{sec:dlog_candidates}

\paragraph{Tadpole families $R_1,R_2$ and $d\log$ candidates}
The master integrals for the $\tau$-dependent subfamily are denoted by $I^{(R)}_{\bm n}$, where the index vector $\bm n=\{n_3,n_4\}$ labels the Hankel function orders.

The tadpole terms $R_1,R_2$ can be treated similarly at the integrand level, by isolating the one-time ``tree-level-like'' block
\begin{align}
	I^{(R)}_{\bm n}(x_2;P_0)\equiv \int_{-\infty}^0 \di\tau\,\hat{I}^{(R)}_{\bm n},
\end{align}
where
\begin{align}
	\hat{I}^{(R)}_{\bm n} \equiv e^{-\ii P_0\tau}\,(-\tau)^{a_0}\,h_{n_3}^{(1)}(-x_2\tau)h_{n_4}^{(2)}(-x_2\tau).
\end{align}
For fixed $(\nu,a_0)$, the $I^{(R)}_{\bm n}$ form a finite basis and satisfy a first-order system
\begin{align}
	\di I^{(R)}_{\bm m}=\left(\di\Omega^{(R)}_{\bm m\bm n}\right)\,I^{(R)}_{\bm n},
\end{align}
where $\di\Omega^{(R)}$ is built from $\di\log P_0$, $\di\log x_2$ and $\di\log(P_0\pm 2x_2)$.
The differential matrix decomposes as:
\begin{align}
	\di\Omega^{(R)}_{\bm m\bm n} = \left(\Omega_{P_0}\right)_{\bm m\bm n} \di\log P_0 + \left(\Omega_{x_2}\right)_{\bm m\bm n} \di\log x_2 + \left(\Omega_{-}\right)_{\bm m\bm n} \di\log\left(P_0-2x_2\right) + \left(\Omega_{+}\right)_{\bm m\bm n} \di\log\left(P_0+2x_2\right),
\end{align}
with coefficient matrices:
\begin{align}
	\Omega_{P_0} &= \frac{-a_0+2\nu}{2} 
	- \frac{1+2\nu}{4} (\sigma_3 \otimes 1 + 1 \otimes \sigma_3) \n\\
	&\quad + \frac{i(1+2\nu)}{4} (\sigma_1 \otimes \sigma_2 + \sigma_2 \otimes \sigma_1)
	+ \frac{a_0-2\nu}{2} \sigma_2 \otimes \sigma_2, \\
	\Omega_{x_2} &= -(1+2\nu) + \frac{1+2\nu}{2}(\sigma_3\otimes 1 + 1\otimes \sigma_3), \\
	\Omega_{+} &= \frac{-a_0+2\nu}{4}(1-\sigma_2)\otimes(1-\sigma_2) \n\\
	&\quad + \frac{1+2\nu}{8}\left[ (i\sigma_1-\sigma_3)\otimes(1-\sigma_2) + (1-\sigma_2)\otimes(i\sigma_1-\sigma_3) \right], \\
	\Omega_{-} &= \frac{-a_0+2\nu}{4}(1+\sigma_2)\otimes(1+\sigma_2) \n\\
	&\quad - \frac{1+2\nu}{8}\left[ (i\sigma_1+\sigma_3)\otimes(1+\sigma_2) + (1+\sigma_2)\otimes(i\sigma_1+\sigma_3) \right].
\end{align}

Furthermore, to study the asymptotic behavior at spatial infinity $x_2 \to \infty$, we can perform the variable transformation $x_2 = 1/t$. By doing so, we can extract the residue matrix $\Omega_\infty$ associated with the $t \to 0$ limit, which is defined as the coefficient of $\di\log t$. Substituting $x_2 = 1/t$ into the relevant differential forms yields:
\begin{align}
	\di\log x_2 &= -\di\log t, \n\\
	\di\log(P_0 \pm 2x_2) &= \di\log\left(\frac{P_0 t \pm 2}{t}\right) = \di\log\left(t \pm \frac{2}{P_0}\right) - \di\log t.
\end{align}
In the limit $t \to 0$, the $t \pm 2/P_0$ terms approach constants and do not contribute to the $\di\log t$ singularity. Thus, collecting the coefficients of $-\di\log t$ from $\di\Omega^{(R)}$, we obtain the definition of $\Omega_\infty$ as a simple linear combination:
\begin{align}
	\Omega_\infty = - \left( \Omega_{x_2} + \Omega_{+} + \Omega_{-} \right).
\end{align}
This matrix $\Omega_\infty$ governs the simple pole at $x_2 \to \infty$.

The basis is formed by tensor products of Pauli matrices, where $1$ denotes the $2\times 2$ identity matrix.
The matrix $\Omega_{x_2}$ is consistent with the integrability condition $d\Omega - \Omega \wedge \Omega = 0$.
Note that $\partial_{P_0}\Omega$ contains a pole at $P_0=0$ from $\Omega_{P_0}$. To construct a valid $d\log$ candidate for $x_2$ integration, one must subtract this contribution:
\begin{align}
	\left(\partial_{P_0}\Omega^{(R)}_{\bm m\bm n} - \frac{1}{P_0}(\Omega_{P_0})_{\bm m\bm n}\right) \di x_2 \sim \left( \tau\delta_{\bm m\bm n} - \frac{\ii}{P_0}(\Omega_{P_0})_{\bm m\bm n}\right) \di x_2.
\end{align}
The combined twist can be packaged as
\begin{align}
	\nu^{(R)}_{\bm n} = x_1^{-b_0}x_2^{-b_0}\,\mG^{-\ep}\mK^{\ep}\,I^{(R)}_{\bm n}.
\end{align}
Some candidate $d\log$-type integrands for $R_1$ in $d=3-2\ep$ are (up to overall constants)
\begin{align}
&	k_s\,\hat{R}_1[\{\bm n\},\{0\},\{2,2\};d=3] = \frac{\di x_1\,\di x_2}{x_1 x_2}\,x_1^{-b_0}x_2^{-b_0}\,I^{(R)}_{\bm n}\,\mG^{-\ep}\mK^{\ep}, \n\\
-	& k_s\,\hat{R}_1[\{\bm n\},\{1\},\{2,1\};d=3] -   \ii	\frac{k_s}{P_0} (\Omega_{P_0})_{\bm n\bm m} \hat{R}_1[\{\bm m\},\{0\},\{2,1\};d=3]     \n\\
& ~~~~~~~~~~~~~= \frac{\di x_1}{x_1}\,\left(\tau\delta_{\bm n\bm m} - \frac{\ii}{P_0}(\Omega_{P_0})_{\bm n\bm m}\right)\,\di x_2 \,x_1^{-b_0}x_2^{-b_0}\,I^{(R)}_{\bm m}\,\mG^{-\ep}\mK^{\ep}.
\end{align}

Besides the basic candidates above, one can still generate further $d\log$-type combinations from the $t=1/x_2$ representation.

\textbf{A further candidate from the $x_2$ equation.} \\
Using $x_2=1/t$, the $x_2$ part of the differential system becomes
\begin{align}
	(\partial_{x_2} I^{(R)}_{\bm m})\,\di x_2
	&= \left(
	-\Omega_\infty \frac{\di x_2}{x_2}
	+ \Omega_+ \di\log\left(t+\frac{2}{P_0}\right)
	+ \Omega_- \di\log\left(t-\frac{2}{P_0}\right)
	\right)_{\bm m\bm n} I^{(R)}_{\bm n}.
\end{align}
\textbf{A further construction starting from $\hat{R}_1[\{\bm n\},\{0\},\{2,1\};d=3]$.} \\
Now start instead from
\begin{align}
	k_s\,\hat{R}_1[\{\bm n\},\{0\},\{2,1\};d=3]
	&= -\frac{\di x_1\,\di t}{x_1\,t^2}\,
	x_1^{-b_0}t^{b_0}\,I^{(R)}_{\bm n}\,\mG^{-\ep}\mK^{\ep}.
\end{align}
Performing IBP once more on the $t$ variable gives
\begin{align}
	&(b_0-1)\,k_s\,\hat{R}_1[\{\bm n\},\{0\},\{2,1\};d=3]
	\sim \frac{\di x_1\,\di t}{x_1\,t}\,
	x_1^{-b_0}t^{b_0}\,\partial_t I^{(R)}_{\bm n}\,\mG^{-\ep}\mK^{\ep} \n\\
	&= \frac{\di x_1\,\di t}{x_1}\,x_1^{-b_0}t^{b_0}
	\Bigg[
	(\Omega_\infty)_{\bm n\bm m}\frac{1}{t^2}
	+(\Omega_+)_{\bm n\bm m}\frac{1}{t\left(t+\frac{2}{P_0}\right)}
	+(\Omega_-)_{\bm n\bm m}\frac{1}{t\left(t-\frac{2}{P_0}\right)}
	\Bigg]
	I^{(R)}_{\bm m}\,\mG^{-\ep}\mK^{\ep}.
\end{align}
Here the first term is the non-$d\log$ obstruction, proportional to $\Omega_\infty\,\di t/t^2$. By contrast,
\begin{align}
	\frac{\di t}{t\left(t\pm \frac{2}{P_0}\right)}
	&= \pm \frac{P_0}{2}
	\left(
	\frac{\di t}{t}
	-\frac{\di t}{t\pm \frac{2}{P_0}}
	\right) \n\\
	&= \pm \frac{P_0}{2}
	\left[
	\di\log t-\di\log\left(t\pm \frac{2}{P_0}\right)
	\right],
\end{align}
or equivalently
\begin{align}
	\frac{2}{P_0}\,\frac{\di t}{t\left(t+\frac{2}{P_0}\right)}
	&= \di\log t-\di\log\left(t+\frac{2}{P_0}\right), \n\\
	\frac{2}{P_0}\,\frac{\di t}{t\left(t-\frac{2}{P_0}\right)}
	&= -\di\log t+\di\log\left(t-\frac{2}{P_0}\right).
\end{align}
Thus, after subtracting the $\Omega_\infty\,\di t/t^2$ piece and absorbing the overall normalization, one gets another valid $d\log$ construction:
\begin{align}
	&-\ii\,\frac{2k_s}{P_0}\left(
	(b_0-1)\delta_{\bm n\bm m}
	-(\Omega_\infty)_{\bm n\bm m}
	\right)\hat{R}_1[\{\bm m\},\{0\},\{2,1\};d=3] \n\\
	&= -\ii\,\frac{\di x_1\,\di t}{x_1}\,x_1^{-b_0}t^{b_0}
	\Bigg[
	(\Omega_+)_{\bm n\bm m}\frac{2/P_0}{t\left(t+\frac{2}{P_0}\right)}  
	+(\Omega_-)_{\bm n\bm m}\frac{2/P_0}{t\left(t-\frac{2}{P_0}\right)}
	\Bigg]
	I^{(R)}_{\bm m}\,\mG^{-\ep}\mK^{\ep}.
\end{align}
The two kernels on the right-hand side are now directly normalized
as linear combinations of $\di\log t$ and $\di\log(t\pm 2/P_0)$,
with coefficients independent of the external kinematics.

In $d=1-2\ep$, one can also use the Baikov building blocks
$\frac{k_s\sqrt{z_1}\,\di z_2}{\mG}$ and $\frac{k_s\sqrt{z_2}\,\di z_1}{\mG}$ to form additional $d\log$ candidates such as
\begin{align}
	&\hat{R}_1[\{\bm n\},\{0\},\{1,0\};d=1]
	= \frac{\di z_1}{z_1}\,\frac{k_s\sqrt{z_1}\,\di z_2}{\mG}\,
	z_1^{-b_0/2}z_2^{-b_0/2}\,I^{(R)}_{\bm n}\,\mG^{-\ep}\mK^{\ep}, \n \\
	&\hat{R}_1[\{\bm n\},\{0\},\{0,1\};d=1]
	= \frac{\di z_2}{z_2}\,\frac{k_s\sqrt{z_2}\,\di z_1}{\mG}\,
	z_1^{-b_0/2}z_2^{-b_0/2}\,I^{(R)}_{\bm n}\,\mG^{-\ep}\mK^{\ep},  \n \\
	-&\hat{R}_1[\{\bm{n}\},\{1\},\{0,0\};d=1]-  \frac{\ii}{P_0} (\Omega_{P_0})_{\bm n\bm m} \hat{R}_1[\{\bm{n}\},\{0\},\{0,0\};d=1]   \n\\                       &~~~~~~~~~~~~~~~~~~~~~= \left(\tau\delta_{\bm n\bm m} - 	 \frac{\ii}{P_0}(\Omega_{P_0})_{\bm n\bm m}\right)\,\di x_2\,\frac{k_s\sqrt{z_2}\,\di z_1}{\mG}\,
	z_1^{-b_0/2}z_2^{-b_0/2}\,I^{(R)}_{\bm m}\,\mG^{-\ep}\mK^{\ep}.
\end{align}


\section{Differential equations for the tadpole master integrals}
\label{sec:deR1}
\subsection{Master-integral definitions}
The Kira implementation uses shifted indices
\begin{align}
	(-\tau)^{a_0+a},\qquad |q|^{-b_1},\qquad |\bm q+\bm k_s|^{-b_2},
\end{align}
with shifts (consistent with the conventions in Section~\ref{sec:dlog})
\begin{align}
	b_{10}=0,\qquad b_{20}=-2\nu,\qquad D=3-2\ep.
\end{align}
Note on $d\log$ indices: The R-function representation uses $x$-variables ($z=x^2$). The $d\log$ condition $dz/z = 2 dx/x$ implies an integrand behavior of $1/x^2$ (index 2) in the R-representation, compared to $1/z$ (index 1) in the Baikov representation. Thus, $d\log$ basis elements in this setup correspond to indices $b_i=2$, not 1.

In this convention the tadpole family is
\begin{align}
	R_1[\{n_3,n_4\},\{a\},\{b_1,b_2\}] \equiv &~e^{\pi \Im [\nu]}\frac{4i}{\pi}\int_{-\infty}^0 \di\tau \int \FR{\di^D \mb q}{(2\pi)^D}\, e^{-\ii P_0\tau}\nonumber\\
	&\times
	\frac{(-\tau)^{a_0+a}}{|q|^{b_{10}+b_1}\,|\bm q+\bm k_s|^{b_{20}+b_2}}
	h_{n_3}^{(1)}(-|\bm q+\bm k_s|\tau)h_{n_4}^{(2)}(-|\bm q+\bm k_s|\tau).
\end{align}


% 文档结束
\end{document}
